
\documentclass[10pt, a4paper, roman, french]{moderncv}
\moderncvstyle{classic}                             
\moderncvcolor{purple}                              
\usepackage[utf8]{inputenc}
\usepackage[light]{CormorantGaramond}
\usepackage[T1]{fontenc}
\usepackage[scale=0.75,a4paper]{geometry}
\usepackage{babel}
\usepackage{geometry}
\usepackage{tikz}
\geometry{hmargin=2.5cm,vmargin=1.5cm}

\definecolor{white}{HTML}{DDDDDD}
\definecolor{gray}{HTML}{404040}
\definecolor{purple}{HTML}{8054cc}

\DeclareRobustCommand{\skills}[1]{ 
    \texorpdfstring{\protect\tikz[baseline]{
		\filldraw[fill=white, draw=gray] (0,0) rectangle (5,0.175);
		\draw[fill=purple](0,0) rectangle ({#1},0.175);
    }}{skill level {#1}}
}


%----------------------------------------------------------------------------------
%  description : string;
%  company : string;  
%  location : string;
%  date : string; 
%} 
%----------------------------------------------------------------------------------
\firstname{Émile}
\familyname{Trotignon}
\mobile{+33 7 82 89 83 58}                          
\extrainfo{Born 30 of July 1999}
\email{emile.trotignon@gmail.com}                                
% \photo[64pt][0.4pt]{Image} % 
\begin{document}
	\makecvtitle
	\section{Formation}
	
		\cventry{2019 -- 2020 }{Computer Science BS}{École Normale Supérieure Paris-Saclay}{}{}{}
	
		\cventry{2018 -- 2019 }{Second year of Computer Science and Mathematics BS}{Université Lyon 1 Claude-Bernard}{}{}{}
	
		\cventry{2017 -- 2018 }{First year of engineering BS}{Jean-Perrin preparatory school}{}{}{}
	
		\cventry{2016 -- 2017 }{High school diploma with science focus}{Lycée La Trinité}{}{}{}
	
	\section{Experience}
	
		\cventry{Summer 2020}{Research internship in computational geometry }{LIRIS laboratory}{Lyon, France}{}{Six weeks internship tutored by David Coeurjolly and Vincent Nivoliers.
                          My goal during this internship was to uniformly sample the surface of a mesh that could have defects.
                          I did a lot of C++ programming during those six weeks, and I used tools such as Polyscope and LIBIGL.
                          My internship report is available here : \url{/files/rapport.pdf}}
	 
		\cventry{Marsh 2020}{Fullstack Node.js developper }{Junior entreprise of ENS Paris-Saclay}{}{}{During a six week mission for the junior entreprise of ENS Paris-Saclay, 
                         I contributed to the developpement of the website of Expert People, a new freelancing plateforme.
                         The technology used were Node.js and Express.js.
                         For instance I programmed a way to fill out the resume of an expert in one click by using his downloadable Linkedin resume in PDF format.
                         Expert People's website (in french) : \url{https://expertpeople.co/}}
	 
		\cventry{January 2020}{ICPC SWERC 2019-2020 }{}{Télécom Paris}{}{University programmation/algorithms competition.
                             Participation in teams of three student.
                             Ranked 37 on 95 teams representing universities from multiple european countries.}
	 
		\cventry{Summer 2019}{Intern C\# developper }{Eternix Ldt.}{Tel Aviv, Israel}{}{Two month internship. HSLS shaders, introduction to DirectX and OpenCV, Windows Form developpement.
                             Greatly rewarding experience abroad.}
	 
		\cventry{July 2018}{Front end developer }{École Nationale Supérieure des Sciences de l'Information et des Bibliothèques}{Lyon, France}{}{ I contributed to the graphical intergration of the new website of the head librarian school for a month. You can see their website here : \url{http://www.enssib.fr/}}
	 

	\section{Languages}
		
	    \cvitem{English}{Fluent}
		
	    \cvitem{French}{Native}
			
	\section{Technical skills} 
	
		\subsection{C++\hfill\skills{4)}}
			I followed a C++ course at university and I also used C++ both for competitive programmation and personnal projects. 
                           
                           I also know C ANSI.
	
		\subsection{Python\hfill\skills{2)}}
			I started Python as a hobby in 2013. I also followed courses on scientific computing.
                       
                       I have a few Python projects available on my GitHub page : 

                       \url{https://github.com/EmileTrotignon?tab=repositories&q=&type=&language=python}
	
		\subsection{C\#\hfill\skills{3)}}
			Good knowledge of C\# : professional experience of 2 month.

                         Two learning projects made in this language :
                         \url{https://github.com/EmileTrotignon/InterpolationShaders}

This project displays a simple texture with DirectX and implement Lanczos interpolation in pixel shaders.

\url{https://github.com/EmileTrotignon/CSGoban}

Small WinForm application allowing to play the game of Go.
	
		\subsection{OCaml\hfill\skills{4)}}
			I had an OCaml course in my first year. 
                
                         In my third year I programmed a compiler for a subset of the C language :

                         \url{https://github.com/EmileTrotignon/mcc}
                       I also published a package on Opam, the Ocaml package manager : \url{https://github.com/EmileTrotignon/ocaml_template_engine}
	
		\subsection{Web development\hfill\skills{3)}}
			Front-end : good knowledge of HTML and CSS : one month experience during summer 2018

                         Back-end : Professional experience developping a Node.js web app
	
		\subsection{Unix\hfill\skills{2)}}
			I know how to use a Unix system with the command line : file manipulation, Git, SSH
	
		\subsection{Scala\hfill\skills{2)}}
			I developped a tower-defense game in Scala during a team project at school. \url{https://github.com/EmileTrotignon/emiles-s-tower-defense}
	
\end{document}